\subsection{Media}

A medium is any material that can exist, naturally or artificially. Air,
water, copper, iron, wood, helium, and vacuum itself are all different media.
In relation to electromagnetics, a medium is characterized by how it responds
to an electric or magnetic function. An applied field induces a response in
the medium's constituent charges (electric dipoles, magnetic dipoles) and
this response in turn modifies the total electric and magnetic flux density
functions. The two functions describing this response are the \textbf{polarization
function} \( \mathbf{P(x,y,z,t)} \), measured in \( \frac{C}{m^2} \), and
the \textbf{magnetization function} \( \mathbf{M(x,y,z,t)} \), measured in
\( \frac{A}{m} \).
There are two ways to express the total flux density functions \( \mathbf{D} \) and
\( \mathbf{B} \):
\begin{itemize}
  \item \textbf{Fundamental Form}: Polarization and Magnetization are explicitly
    included in the equations:
    \begin{align*}
      \mathbf{D(x,y,z,t)} &= \epsilon_0 \mathbf{E(x,y,z,t)} + \mathbf{P(x,y,z,t)} \tag{10a} \\
      \mathbf{B(x,y,z,t)} &= \mu_0 \left[\mathbf{H(x,y,z,t)} + \mathbf{M(x,y,z,t)}\right] \tag{10b}
    \end{align*}
    Here \( \epsilon_0 \approx 8.854 \times 10^{-12} \, \frac{F}{m} \) is known
    as the \textbf{permittivity of free space} and \( \mu_0 \approx 4\pi \times 10^{-7} \, \frac{H}{m} \)
    is known as the \textbf{permeability of free space}. These are fixed, universal
    constants describing vacuum.
  \item \textbf{Tensor Form}: Polarization and Magnetization are absorbed into the
    permittivity and permeability tensors:
    \begin{align*}
      \mathbf{D(x,y,z,t)} &= \mathbf{K}_E(x,y,z,t) \cdot \mathbf{E(x,y,z,t)} \tag{11a} \\
      \mathbf{B(x,y,z,t)} &= \mathbf{K}_M(x,y,z,t) \cdot \mathbf{H(x,y,z,t)} \tag{11b}
    \end{align*}
    (\(\cdot\) here is the matrix times vector product). This instantaneous product
    is exact only for a \textbf{non-dispersive} medium (below). A dispersive
    medium responds to the past of the field, so in time the product becomes
    a convolution and in frequency it becomes a product again:
    \begin{align*}
      \mathbf{D}(t) &= \int_0^{\infty} \mathbf{K}_E(t') \cdot \mathbf{E}(t - t')\, dt' \hspace{1cm}
      \mathbf{D}(\omega) = \mathbf{K}_E(\omega) \cdot \mathbf{E}(\omega)
    \end{align*}
    \( \mathbf{K}_E \) (permittivity) and \( \mathbf{K}_M \)
    (permeability) are, in general, \textbf{rank-2 tensors} i.e. \( 3 \times 3 \)
    matrices with 9 elements each:
    \begin{align*}
      \mathbf{K}_E(x,y,z,t) = \begin{bmatrix}
        K_{E,xx} & K_{E,xy} & K_{E,xz} \\
        K_{E,yx} & K_{E,yy} & K_{E,yz} \\
        K_{E,zx} & K_{E,zy} & K_{E,zz}
      \end{bmatrix}, \qquad
      \mathbf{K}_M(x,y,z,t) = \begin{bmatrix}
        K_{M,xx} & K_{M,xy} & K_{M,xz} \\
        K_{M,yx} & K_{M,yy} & K_{M,yz} \\
        K_{M,zx} & K_{M,zy} & K_{M,zz}
      \end{bmatrix}
    \end{align*}
    so that, written out fully with no shorthand,
    \begin{align*}
      \begin{bmatrix}
        D_x(t) \\
        D_y(t) \\
        D_z(t)
      \end{bmatrix}
      &=
      \begin{bmatrix}
        K_{E,xx} & K_{E,xy} & K_{E,xz} \\
        K_{E,yx} & K_{E,yy} & K_{E,yz} \\
        K_{E,zx} & K_{E,zy} & K_{E,zz}
      \end{bmatrix}
      \begin{bmatrix}
        E_x(t) \\
        E_y(t) \\
        E_z(t)
      \end{bmatrix} \\
      \begin{bmatrix}
        D_x(t) \\
        D_y(t) \\
        D_z(t)
      \end{bmatrix}
      &=
      \begin{bmatrix}
        K_{E,xx} E_x(t) + K_{E,xy} E_y(t) + K_{E,xz} E_z(t) \\
        K_{E,yx} E_x(t) + K_{E,yy} E_y(t) + K_{E,yz} E_z(t) \\
        K_{E,zx} E_x(t) + K_{E,zy} E_y(t) + K_{E,zz} E_z(t)
      \end{bmatrix} \\
      \begin{bmatrix}
        B_x(t) \\
        B_y(t) \\
        B_z(t)
      \end{bmatrix}
      &=
      \begin{bmatrix}
        K_{M,xx} & K_{M,xy} & K_{M,xz} \\
        K_{M,yx} & K_{M,yy} & K_{M,yz} \\
        K_{M,zx} & K_{M,zy} & K_{M,zz}
      \end{bmatrix}
      \begin{bmatrix}
        H_x(t) \\
        H_y(t) \\
        H_z(t)
      \end{bmatrix} \\
        \begin{bmatrix}
        B_x(t) \\
        B_y(t) \\
        B_z(t)
      \end{bmatrix}
      &=
      \begin{bmatrix}
        K_{M,xx} H_x(t) + K_{M,xy} H_y(t) + K_{M,xz} H_z(t) \\
        K_{M,yx} H_x(t) + K_{M,yy} H_y(t) + K_{M,yz} H_z(t) \\
        K_{M,zx} H_x(t) + K_{M,zy} H_y(t) + K_{M,zz} H_z(t)
      \end{bmatrix}
    \end{align*}
    Notice that in general, \( D_x \) can depend on \emph{all three} components of
    \( \mathbf{E} \), not just \( E_x \) and similarly \( B_x \) can depend on all three
    components of \( \mathbf{H} \), not just \( H_x \). This is the mathematical statement
    that \( \mathbf{D} \) and \( \mathbf{E} \) or \( \mathbf{B} \) and \( \mathbf{H} \)
    need not point in the same direction.
\end{itemize}

\subsubsection{Classifying a Medium}

Every medium's \( \mathbf{K}_E \) and \( \mathbf{K}_M \) can be
classified along (at least) five independent axes, each describing what the
9 tensor elements are allowed to depend on:
\begin{itemize}
  \item \textbf{Homogeneity} [dependence on position]: A medium is
    \textbf{homogeneous} if no element of \( \mathbf{K}_E \) or
    \( \mathbf{K}_M \) depends on \( (x,y,z) \). It is
    \textbf{inhomogeneous} if the elements are functions of position,
    \( K_{E,ij}(x,y,z) \). A single crystal of pure quartz, block
    of solid copper or air inside a room are examples of homogeneous media.
    Metamaterials like Graded Index (GRIN) lenses, Earth's Atmosphere
    (Ionosphere), Biological tissue are examples of inhomogeneous media.
  \item \textbf{Isotropy} [dependence on direction]: A medium is
    \textbf{isotropic} if its tensor is diagonal with all three diagonal
    elements equal \( K_{E,ij} = K_E \, I_{ij} \) (where
    \( I_{ij} \) is the Kronecker delta: 1 if \( i=j \), 0 otherwise or
    identity in matrix terms). This collapses the tensor to a single scalar
    \( K_E \), and \( \mathbf{D} \) and \( \mathbf{E} \) become
    collinear. A medium is \textbf{anisotropic} if the diagonal elements
    are unequal and/or any off-diagonal element is nonzero. Glass, Water
    and Air are Isotropic media. Calcite and Quartz Crystal are examples
    of Anisotropic media.
  \item \textbf{Linearity} [dependence on field magnitude]: A medium is
    \textbf{linear} if \( \mathbf{K}_E \) and \( \mathbf{K}_M \)
    do not depend on the magnitude of \( \mathbf{E} \) or \( \mathbf{H} \)
    themselves. It is \textbf{nonlinear} if they do, \( K_{E,ij}(\mathbf{E}) \).
    All ordinary dielectrics are mostly linear. Kerr Media, Second-Harmonic Generation
    Crystals and Ferromagnetic cores are examples of Nonlinear media.
  \item \textbf{Dispersion} [dependence on frequency]: A medium is
    \textbf{non-dispersive} if \( \mathbf{K}_E \) and
    \( \mathbf{K}_M \) do not depend on frequency \( \omega \). It is
    \textbf{dispersive} if they do, \( K_{E,ij}(\omega) \). Equivalently in
    the time domain, the medium's response depends on the field's past
    values (memory), not just its instantaneous value. Vacuum is exactly
    non-dispersive; air is very nearly so. Glass, Water (in Microwave and
    Optical frequencies) and Metals are examples of Dispersive media.
  \item \textbf{Time-Variance} [dependence on time]: A medium is
    \textbf{time-invariant} if \( \mathbf{K}_E \) and
    \( \mathbf{K}_M \) do not explicitly depend on \( t \). It is
    \textbf{time-varying} if the medium's own properties are being actively
    modulated in time \( K_{E,ij}(t) \), distinct from dispersion,
    which is memory of the field's past, not modulation of the medium
    itself. A material left alone (not modulated) is time-invariant. A
    varactor whose bias is modulated, a plasma with time varying electron
    density and space-time metamaterials are examples of Time-Varying media.
\end{itemize}
\textbf{These axes are independent of each other}. The most general case is when a material
is inhomogeneous, anisotropic, nonlinear, dispersive, and time-varying, with all 9
tensor elements independent (anisotropy) and each an arbitrary function of position,
field, frequency, and time. The simplest case is homogeneous, isotropic, linear,
non-dispersive, and time-invariant: the three diagonal elements are one equal
constant and the six off-diagonal elements are zero.
Most real media fall somewhere in between these two extremes.

\subsubsection{Vacuum as the Isotropic, Homogeneous Reference}

Vacuum is chosen as the reference medium. It is homogeneous, isotropic,
linear, non-dispersive, and time-invariant, so its tensor collapses
to the constant, diagonal, equal-entry case:
\begin{align*}
  \mathbf{K}_{E,\text{vacuum}} = \epsilon_0 \, \mathbb{I} = \begin{bmatrix}
    \epsilon_0 & 0 & 0 \\ 0 & \epsilon_0 & 0 \\ 0 & 0 & \epsilon_0
  \end{bmatrix}, \qquad
  \mathbf{K}_{M,\text{vacuum}} = \mu_0 \, \mathbb{I} = \begin{bmatrix}
    \mu_0 & 0 & 0 \\ 0 & \mu_0 & 0 \\ 0 & 0 & \mu_0
  \end{bmatrix}
\end{align*}
Since this tensor is proportional to the identity matrix \( \mathbb{I} \),
writing it as the bare scalar \( \epsilon_0 \) (or \( \mu_0 \)) loses no
information. Any other medium's tensor is compared to this reference by
factoring it out:
\begin{align*}
  \mathbf{K}_E(x,y,z,t) &= \epsilon_0 \, \mathbf{K}_{ER}(x,y,z,t) \\
  \mathbf{K}_M(x,y,z,t) &= \mu_0 \, \mathbf{K}_{MR}(x,y,z,t)
\end{align*}
where \( \mathbf{K}_{ER} \) and \( \mathbf{K}_{MR} \) are the
dimensionless \textbf{relative permittivity} and \textbf{relative
permeability} tensors, subject to the same classification above. Vacuum is
recovered exactly when \( \mathbf{K}_{ER} = \mathbf{K}_{MR} =
\mathbb{I} \).

\subsubsection{Recovering P and M}
Both forms must be equal.
\begin{itemize}
  \item Solving for \( \mathbf{P} \): Equation (10a)
    and Equation (11a) must be equal.
    \begin{align*}
      \epsilon_0
      \begin{bmatrix}
        E_x(t) \\
        E_y(t) \\
        E_z(t)
      \end{bmatrix} +
      \begin{bmatrix}
        P_x(t) \\
        P_y(t) \\
        P_z(t)
      \end{bmatrix}
      &= \epsilon_0
      \begin{bmatrix}
        (K_{ER})_{xx} & (K_{ER})_{xy} & (K_{ER})_{xz} \\
        (K_{ER})_{yx} & (K_{ER})_{yy} & (K_{ER})_{yz} \\
        (K_{ER})_{zx} & (K_{ER})_{zy} & (K_{ER})_{zz}
      \end{bmatrix}
      \begin{bmatrix}
        E_x(t) \\
        E_y(t) \\
        E_z(t)
      \end{bmatrix} \\
      \begin{bmatrix}
        E_x(t) + \frac{P_x(t)}{\epsilon_0} \\
        E_y(t) + \frac{P_y(t)}{\epsilon_0} \\
        E_z(t) + \frac{P_z(t)}{\epsilon_0}
      \end{bmatrix}
      &=
      \begin{bmatrix}
        (K_{ER})_{xx} E_x(t) + (K_{ER})_{xy} E_y(t) + (K_{ER})_{xz} E_z(t) \\
        (K_{ER})_{yx} E_x(t) + (K_{ER})_{yy} E_y(t) + (K_{ER})_{yz} E_z(t) \\
        (K_{ER})_{zx} E_x(t) + (K_{ER})_{zy} E_y(t) + (K_{ER})_{zz} E_z(t)
      \end{bmatrix} \hspace{1cm} \left[\div \epsilon_0\right] \\
      \frac{1}{\epsilon_0}\begin{bmatrix}
        P_x(t) \\
        P_y(t) \\
        P_z(t)
      \end{bmatrix}
      &=
      \begin{bmatrix}
        \left((K_{ER})_{xx}E_x(t) + (K_{ER})_{xy}E_y(t) + (K_{ER})_{xz}E_z(t)\right) - E_x(t) \\
        \left((K_{ER})_{yx}E_x(t) + (K_{ER})_{yy}E_y(t) + (K_{ER})_{yz}E_z(t)\right) - E_y(t) \\
        \left((K_{ER})_{zx}E_x(t) + (K_{ER})_{zy}E_y(t) + (K_{ER})_{zz}E_z(t)\right) - E_z(t)
      \end{bmatrix} \\
      \frac{1}{\epsilon_0}\begin{bmatrix}
        P_x(t) \\
        P_y(t) \\
        P_z(t)
      \end{bmatrix}
      &=
      \begin{bmatrix}
        \left((K_{ER})_{xx}E_x(t) - E_x(t)\right) + (K_{ER})_{xy}E_y(t) + (K_{ER})_{xz}E_z(t) \\
        (K_{ER})_{yx}E_x(t) + \left((K_{ER})_{yy}E_y(t) - E_y(t)\right) + (K_{ER})_{yz}E_z(t) \\
        (K_{ER})_{zx}E_x(t) + (K_{ER})_{zy}E_y(t) + \left((K_{ER})_{zz}E_z(t) - E_z(t)\right)
      \end{bmatrix} \\
      \frac{1}{\epsilon_0}\begin{bmatrix}
        P_x(t) \\
        P_y(t) \\
        P_z(t)
      \end{bmatrix}
      &=
      \begin{bmatrix}
        \left((K_{ER})_{xx} - 1\right)E_x(t) + (K_{ER})_{xy}E_y(t) + (K_{ER})_{xz}E_z(t) \\
        (K_{ER})_{yx}E_x(t) + \left((K_{ER})_{yy} - 1 \right)E_y(t) + (K_{ER})_{yz}E_z(t) \\
        (K_{ER})_{zx}E_x(t) + (K_{ER})_{zy}E_y(t) + \left((K_{ER})_{zz} - 1\right)E_z(t)
      \end{bmatrix} \\
      \frac{1}{\epsilon_0}\begin{bmatrix}
        P_x(t) \\
        P_y(t) \\
        P_z(t)
      \end{bmatrix}
      &=
      \begin{bmatrix}
        (K_{ER})_{xx} - 1 & (K_{ER})_{xy} & (K_{ER})_{xz} \\
        (K_{ER})_{yx} & (K_{ER})_{yy} - 1 & (K_{ER})_{yz} \\
        (K_{ER})_{zx} & (K_{ER})_{zy} & (K_{ER})_{zz} - 1
      \end{bmatrix}
      \begin{bmatrix}
        E_x(t) \\
        E_y(t) \\
        E_z(t)
      \end{bmatrix} \\
      \begin{bmatrix}
        P_x(t) \\
        P_y(t) \\
        P_z(t)
      \end{bmatrix}
      &= \epsilon_0
      \begin{bmatrix}
        (K_{ER})_{xx} - 1 & (K_{ER})_{xy} & (K_{ER})_{xz} \\
        (K_{ER})_{yx} & (K_{ER})_{yy} - 1 & (K_{ER})_{yz} \\
        (K_{ER})_{zx} & (K_{ER})_{zy} & (K_{ER})_{zz} - 1
      \end{bmatrix}
      \begin{bmatrix}
        E_x(t) \\
        E_y(t) \\
        E_z(t)
      \end{bmatrix} \hspace{1cm} \left[\times \epsilon_0\right]
    \end{align*}
    If you define \( \mathbf{X}_E \equiv \mathbf{K}_{ER} - \mathbb{I} \) as the
  \textbf{electric susceptibility tensor}, then the above can be written more compactly as
    \begin{align*}
      \begin{bmatrix}
        P_x(t) \\
        P_y(t) \\
        P_z(t)
      \end{bmatrix}
      &= \epsilon_0
      \begin{bmatrix}
        (X_E)_{xx} & (X_E)_{xy} & (X_E)_{xz} \\
        (X_E)_{yx} & (X_E)_{yy} & (X_E)_{yz} \\
        (X_E)_{zx} & (X_E)_{zy} & (X_E)_{zz}
      \end{bmatrix}
      \begin{bmatrix}
        E_x(t) \\
        E_y(t) \\
        E_z(t)
      \end{bmatrix}
    \end{align*}
    where \( \mathbf{X}_E \) is also a rank-2 tensor with 9 elements and
    equals,
    \begin{align*}
      \mathbf{X}_E(x,y,z,t) = \begin{bmatrix}
        (K_{ER})_{xx} - 1 & (K_{ER})_{xy} & (K_{ER})_{xz} \\
        (K_{ER})_{yx} & (K_{ER})_{yy} - 1 & (K_{ER})_{yz} \\
        (K_{ER})_{zx} & (K_{ER})_{zy} & (K_{ER})_{zz} - 1
      \end{bmatrix}
    \end{align*}
  \item Solving for \( \mathbf{M} \):
    \begin{align*}
      \mu_0 \left(
      \begin{bmatrix}
        H_x(t) \\
        H_y(t) \\
        H_z(t)
      \end{bmatrix} +
      \begin{bmatrix}
        M_x(t) \\
        M_y(t) \\
        M_z(t)
      \end{bmatrix} \right)
      &= \mu_0
      \begin{bmatrix}
        (K_{MR})_{xx} & (K_{MR})_{xy} & (K_{MR})_{xz} \\
        (K_{MR})_{yx} & (K_{MR})_{yy} & (K_{MR})_{yz} \\
        (K_{MR})_{zx} & (K_{MR})_{zy} & (K_{MR})_{zz}
      \end{bmatrix}
      \begin{bmatrix}
        H_x(t) \\
        H_y(t) \\
        H_z(t)
      \end{bmatrix} \\
      \begin{bmatrix}
        H_x(t) + M_x(t) \\
        H_y(t) + M_y(t) \\
        H_z(t) + M_z(t)
      \end{bmatrix}
      &=
      \begin{bmatrix}
        (K_{MR})_{xx} H_x(t) + (K_{MR})_{xy} H_y(t) + (K_{MR})_{xz} H_z(t) \\
        (K_{MR})_{yx} H_x(t) + (K_{MR})_{yy} H_y(t) + (K_{MR})_{yz} H_z(t) \\
        (K_{MR})_{zx} H_x(t) + (K_{MR})_{zy} H_y(t) + (K_{MR})_{zz} H_z(t)
      \end{bmatrix} \hspace{1cm} \left[\div \mu_0\right] \\
      \begin{bmatrix}
        M_x(t) \\
        M_y(t) \\
        M_z(t)
      \end{bmatrix}
      &=
      \begin{bmatrix}
          \left((K_{MR})_{xx}H_x(t) + (K_{MR})_{xy}H_y(t) + (K_{MR})_{xz}H_z(t)\right) - H_x(t) \\
          \left((K_{MR})_{yx}H_x(t) + (K_{MR})_{yy}H_y(t) + (K_{MR})_{yz}H_z(t)\right) - H_y(t) \\
          \left((K_{MR})_{zx}H_x(t) + (K_{MR})_{zy}H_y(t) + (K_{MR})_{zz}H_z(t)\right) - H_z(t)
      \end{bmatrix} \\
      \begin{bmatrix}
        M_x(t) \\
        M_y(t) \\
        M_z(t)
      \end{bmatrix}
      &=
      \begin{bmatrix}
          \left((K_{MR})_{xx}H_x(t) - H_x(t) \right) + (K_{MR})_{xy}H_y(t) + (K_{MR})_{xz}H_z(t) \\
          (K_{MR})_{yx}H_x(t) + \left((K_{MR})_{yy}H_y(t) - H_y(t)\right) + (K_{MR})_{yz}H_z(t) \\
          (K_{MR})_{zx}H_x(t) + (K_{MR})_{zy}H_y(t) + \left((K_{MR})_{zz}H_z(t) - H_z(t)\right)
      \end{bmatrix} \\
      \begin{bmatrix}
        M_x(t) \\
        M_y(t) \\
        M_z(t)
      \end{bmatrix}
      &=
      \begin{bmatrix}
          \left((K_{MR})_{xx} - 1\right)H_x(t) + (K_{MR})_{xy}H_y(t) + (K_{MR})_{xz}H_z(t) \\
          (K_{MR})_{yx}H_x(t) + \left((K_{MR})_{yy} - 1\right)H_y(t) + (K_{MR})_{yz}H_z(t) \\
          (K_{MR})_{zx}H_x(t) + (K_{MR})_{zy}H_y(t) + \left((K_{MR})_{zz} - 1\right)H_z(t)
      \end{bmatrix} \\
      \begin{bmatrix}
        M_x(t) \\
        M_y(t) \\
        M_z(t)
      \end{bmatrix}
      &=
      \begin{bmatrix}
        (K_{MR})_{xx} - 1 & (K_{MR})_{xy}
        & (K_{MR})_{xz} \\
        (K_{MR})_{yx} & (K_{MR})_{yy} - 1 & (K_{MR})_{yz} \\
        (K_{MR})_{zx} & (K_{MR})_{zy} & (K_{MR})_{zz} - 1
      \end{bmatrix}
      \begin{bmatrix}
        H_x(t) \\
        H_y(t) \\
        H_z(t)
      \end{bmatrix}
    \end{align*}
    If you define \( \mathbf{X}_M \equiv \mathbf{K}_{MR} - \mathbb{I} \) as the
    \textbf{magnetic susceptibility tensor}, then the above can be written more compactly
    as,
    \begin{align*}
      \begin{bmatrix}
        M_x(t) \\
        M_y(t) \\
        M_z(t)
      \end{bmatrix}
      &=
      \begin{bmatrix}
        (X_M)_{xx} & (X_M)_{xy} & (X_M)_{xz} \\
        (X_M)_{yx} & (X_M)_{yy} & (X_M)_{yz} \\
        (X_M)_{zx} & (X_M)_{zy} & (X_M)_{zz}
      \end{bmatrix}
      \begin{bmatrix}
        H_x(t) \\
        H_y(t) \\
        H_z(t)
      \end{bmatrix}
    \end{align*}
    where \( \mathbf{X}_M \) is also a rank-2 tensor with 9 elements and equals,
    \begin{align*}
      \mathbf{X}_M(x,y,z,t) = \begin{bmatrix}
        (K_{MR})_{xx} - 1 & (K_{MR})_{xy} & (K_{MR})_{xz} \\
        (K_{MR})_{yx} & (K_{MR})_{yy} - 1 & (K_{MR})_{yz} \\
        (K_{MR})_{zx} & (K_{MR})_{zy} & (K_{MR})_{zz} - 1
      \end{bmatrix}
    \end{align*}
\end{itemize}

\paragraph{A Practical Note: Real Media Are Piecewise-Inhomogeneous}
In practice no region of interest is a single, uniform medium. A block of
space spanning open air may contain a concrete building, a glass window,
and a person standing in a doorway. Thus the medium  is \textbf{piecewise-homogeneous}
that is uniform within each material but genuinely a function of position,
\( \mathbf{K}_E(x,y,z) \), \( \mathbf{K}_M(x,y,z) \), over the
region as a whole typically discontinuous at material boundaries. This is
the normal case for real propagation problems (indoor RF coverage, urban
propagation, biomedical imaging) and it is why numerical solvers that time-march
the fields on a spatial grid (say via the central-difference update
equations of the FDTD method) assign \( \mathbf{K}_E \) and
\( \mathbf{K}_M \) \emph{per grid cell} based on which physical material
occupies that cell rather than treating them as global constants pulled outside the
update equation.
